\documentclass[10pt,a4paper]{amsart}
\usepackage[margin=19mm]{geometry}
\usepackage{amsmath,amssymb,mathtools}
\usepackage[scr=rsfs,cal=euler]{mathalfa}
\usepackage{xfrac}
\usepackage[unicode,hidelinks,bookmarks=false]{hyperref}
\newcommand{\ie}{i.e.\ }
\newcommand{\cf}{cf.\ }
\newcommand{\resp}{respectively\ }
\newcommand{\Subsection}{\subsection}
\providecommand{\nobreakdash}{\nobreak\hbox{-}}
\setlength{\emergencystretch}{2em}
\multlinegap0pt
\usepackage{bm}
\usepackage{bbm}
\usepackage{dsfont}
\usepackage{xspace}
\usepackage{cancel}
\usepackage{mathtools}
\usepackage{amsthm}
\makeatletter
\@ifundefined{thm}{\newtheorem{thm}{Thm}}{}
\@ifundefined{lem}{\newtheorem{lem}[thm]{Lemma}}{}
\@ifundefined{prop}{\newtheorem{prop}[thm]{Proposition}}{}
\makeatother
\newcommand{\K}{\mathbb{K}}
\title[Generalized Spectral Bound for Quasi-Twisted Codes]{Generalized Spectral Bound for Quasi-Twisted Codes}
\author{Buket \"Ozkaya}
\date{Source: arXiv:2511.10066v1 (CC BY 4.0)}
\begin{document}
\maketitle

\noindent\emph{Source and licence.} Generalized Spectral Bound for Quasi-Twisted Codes. Version arXiv:2511.10066v1 (CC BY 4.0), distributed under the Creative Commons Attribution 4.0 International licence, as recorded in the arXiv licence metadata for this exact version and shown on the abstract page https://arxiv.org/abs/2511.10066v1. Full rights notice: licenses/arxiv-2511.10066-CC-BY-4.0.txt.\par\smallskip

\noindent\textit{Editorial scope.} Counted unit: the Prop whose source label is \texttt{bound prop} in the article (source0.tex lines 473--486, source label \texttt{bound prop}), delivered together with the article's own complete original proof of it (source0.tex lines 487--526, one \texttt{proof} environment of 40 lines). No mathematics is written, repaired or re-derived: statement and proof are the article's own text line for line. Required context retained uncounted: bound lemma (lines 393--399). Every definition, display and label of those lines is retained unchanged. Omitted from this package and not redistributed: 1--392, 400--472, 527--734. Nothing inside a retained excerpt is omitted or altered: every retained body is delivered exactly as the article's lines are written, so no omission receipt of any kind accompanies this package. Citations: the retained excerpts carry no citation command at all, so no bibliography entry of the article is transcribed and no reference is redistributed. Reference adaptation: none was needed and none was made; every reference inside the delivered unit resolves inside it, so no unresolved reference or citation is produced. Printed slips kept literal: no printed word, symbol, hypothesis or number of the counted statement or of its proof was altered. Layout and class: the article's own class is \texttt{amsart} with packages including float, amsmath, epsfig, latexsym, indentfirst, mathrsfs, amsthm, amssymb, amsfonts, amsbsy, appendix, enumitem, url, multirow; the wrapper is the lane's pinned \texttt{amsart} editorial wrapper at 10pt on A4, which restates the theorem-like environments the retained text needs and adds the article's own macro definitions that the retained text needs (\texttt{\textbackslash K}), with extra wrapper packages (bm, bbm, dsfont, xspace, cancel, mathtools). The delivered numbering is the unit's own. Assets: no retained span contains a figure environment, a graphics-inclusion command, a PDF-page inclusion, a table or a TikZ picture; no image file is copied into this package, the package declares no asset dependency and no image file is redistributed with it. Licence: version arXiv:2511.10066v1 of this article is distributed under the Creative Commons Attribution 4.0 International licence, as recorded in the arXiv licence metadata for this exact version and shown on the abstract page https://arxiv.org/abs/2511.10066v1. A full rights notice is delivered at licenses/arxiv-2511.10066-CC-BY-4.0.txt. Characters: every retained excerpt is pure ASCII, so the delivered engine, XeLaTeX with its default Latin Modern text font, reports no missing character.
\begin{lem}\label{bound lemma}
Let $A$ and $B$ be two matrices with entries from some finite field $\K$.
\begin{itemize}
\item[i.] If $A$ and $B$ have the same number of columns, then the matrix $M={A\choose B}$ satisfies $n_M=\max\{n_A,n_B\}$.
\item[ii.] If we set $M=(A\otimes B)$, then we have $n_M=\min\{n_A,n_B\}$.
\end{itemize}
\end{lem}

\begin{prop}\label{bound prop}
For a positive integer $s$, let $A_1,\ldots, A_s$ be $k_1 \times n_1$ matrices and let $B_1,\ldots, B_s$ be $k_2 \times n_2$ matrices over some finite field $\K$ such that $n_{A_1}\geq n_{A_2}\geq \cdots \geq n_{A_s}$. We define \[M:=\begin{pmatrix}
A_1 \otimes B_1\\
\vdots\\
A_s \otimes B_s
\end{pmatrix}.\]
Then, we have $n_M\geq \min\{n_{A_1},n_{A_2}n_{B_1}, n_{A_3}n_{B_{1+2}},\ldots, n_{A_s}n_{B_{1+2+\cdots +(s-1)}},n_{B_{1+2+\cdots +s}}\},$ where \[B_{1+2+\cdots +j}=\begin{pmatrix}
 B_1\\
 B_2\\
\vdots\\
 B_j
\end{pmatrix},\]
for any $j\in\{1,\ldots,s\}$.
\end{prop}

\begin{proof}
The proof follows by induction on $s$. Notice that the case $s=1$ corresponds to the second part of Lemma \ref{bound lemma}. Now let $s=2$ and $M:=\begin{pmatrix}
A_1 \otimes B_1\\
A_2 \otimes B_2
\end{pmatrix}$ with $n_{A_1}\geq n_{A_2}$. In the cases we will consider below, we use Lemma \ref{bound lemma} and the following fact: if $n_{A_1}\geq n_{B_1}$ and $n_{B_2}\geq n_{A_2}$, then we have $n_M= n_{A_2} n_{B_1}$. To see this, we can take a nonzero vector $c\in\K^{n_1n_2}$ and consider $Mc^\top = 0$. To satisfy this homogeneous equality, all of the $n_1$ blocks of length $n_2$ in $c$ must have weight at least $n_{B_1}+1$ or be zero so that $B_1(c_1)^\top=\cdots=B_1(c_{n_1})^\top=0$ and at least $n_{A_2}+1$ out of these $n_1$ blocks must be nonzero (hence, of weight at least $n_{B_1}+1$) so that $A_2$ is cancelled out. 

\emph{Case 1: When $n_{A_1}\geq n_{A_2}\geq n_{B_{1+2}}$,} the number of linearly independent columns are determined by $B_1$ and $B_2$, which follows by Lemma \ref{bound lemma}. Hence, we get $n_M\geq n_{B_{1+2}}=\max\{n_{B_1},n_{B_2}\}$.

\emph{Case 2: When $n_{B_{1+2}}\geq n_{A_1}\geq n_{A_2}$,} the number of linearly independent columns are determined by $A_1$ and $A_2$, which follows by Lemma \ref{bound lemma}. Here, $n_M\geq\max\{n_{A_1},n_{A_2}\}=n_{A_1}$.

\emph{Case 3: When $n_{A_1}\geq n_{B_{1+2}}\geq n_{A_2}$,} we have to distinguish two possibilities. First, we consider when $n_{B_{1+2}}=n_{B_1}$. This means we either have $n_{A_1}\geq n_{B_1}\geq n_{B_2}\geq n_{A_2}$, implying $n_M= n_{A_2} n_{B_1}$, or $n_{A_1}\geq n_{B_1}\geq n_{A_2}\geq n_{B_2}$, implying $n_M=n_{B_{1+2}}$. Second, when $n_{B_{1+2}}=n_{B_2}$, then we have $n_M= n_{A_2} n_{B_1}$.

Therefore, we obtain $n_M\geq\min\{n_{A_1}, n_{A_2} n_{B_1}, n_{B_{1+2}}\}$ for $s=2$.

Assume that the assertion holds for $s-1$, i.e., 
\begin{equation}\label{ind-hyp}
n_M\geq \min\{n_{A_1},n_{A_2}n_{B_1}, \ldots, n_{A_{s-1}}n_{B_{1+2+\cdots +(s-2)}},n_{B_{1+2+\cdots +(s-1)}}\},
\mbox{\ where\ } M:=\begin{pmatrix}
A_1 \otimes B_1\\
\vdots\\
A_{s-1} \otimes B_{s-1}
\end{pmatrix}.
\end{equation}

By adding $A_s\otimes B_s$ to $M$, we want to show the proposition. Mimicing the case analysis done for $s=2$ above, we consider the following:

\emph{Case 1: When $n_{A_1}\geq \cdots\geq n_{A_S}\geq n_{B_{1+\cdots +s}}$,} the number of linearly independent columns rely on $B_1, \ldots, B_s$, which again comes from Lemma \ref{bound lemma}. Hence, we get $n_M\geq n_{B_{1+\cdots + s}}$.

\emph{Case 2: When $n_{B_{1+\cdots +s}}\geq n_{A_1}\geq \cdots\geq n_{A_s}$,} then we have to consider $A_1,\ldots,A_s$, which again follows by Lemma \ref{bound lemma}. Here, $n_M\geq\max\{n_{A_1},n_{A_2},\ldots, n_{A_s}\}=n_{A_1}$.

\emph{Case 3: When $n_{A_1}\geq \cdots \geq n_{A_{s-1}}\geq n_{B_{1+\cdots +s}}\geq n_{A_s}$,} we have to distinguish two possibilities, as done above. First, we consider when $n_{B_{1+\cdots +s}}=n_{B_{1+\cdots +(s-1)}}$ giving $n_M\geq n_{A_s} n_{B_{1+\cdots +(s-1)}}$ or  $n_M\geq n_{B_{1+\cdots +s}}$. Second, when $n_{B_{1+\cdots +s}}=n_{B_s}$, then we have $n_M\geq n_{A_s} n_{B_{1+\cdots +(s-1)}}$.

The remaining cases correspond to
\begin{itemize}
\item $n_{A_1}\geq \cdots \geq n_{A_{s-2}}\geq n_{B_{1+\cdots +s}}\geq n_{A_{s-1}}\geq n_{A_s}$,\\
 $\vdots$
\item $ n_{A_1}\geq n_{B_{1+\cdots +s}}\geq n_{A_2}\geq \cdots \geq n_{A_s}$,
\end{itemize}
which are covered by the induction hypothesis \eqref{ind-hyp}. More precisely, when $n_{A_1}\geq \cdots \geq n_{A_{s-j-1}}\geq n_{B_{1+\cdots +s}}\geq n_{A_{s-j}}\geq \cdots \geq n_{A_s}$, for some $1\leq j \leq s-2$, we again distinguish two scenarios. If $n_{B_{1+\cdots +s}}=n_{B_{1+\cdots +(s-j-1)}}$, then $n_M\geq n_{A_{s-j}} n_{B_{1+\cdots +(s-j-1)}}$ or  $n_M\geq n_{B_{1+\cdots +s}}$. Otherwise, if $n_{B_{1+\cdots +s}}=n_{B_{s-k}}$, for some $k\in\{0,\ldots, j\}$, then we have $n_M\geq n_{A_{s-j}} n_{B_{1+\cdots +(s-j-1)}}$.
\end{proof}

\end{document}
